\section{The search and its tree}
\label{sec:tree}
The search fills columns from left to right. A node is a partial assignment $p_0,\dots,p_{c-1}$ that is
Costas-valid so far: its rows are distinct, and at every lag the vertical differences are distinct. For
each new column the search tries every unused row, checking the new dot's displacements against a
bitmask of the differences already used at each lag. The first column's row is restricted to
$p_0\le\lfloor(n-1)/2\rfloor$, which uses the vertical-complement symmetry $p\mapsto n-1-p$ to halve the
tree, as in the measurements of Paper~1. Nothing else is pruned.

The implementation, \texttt{estimate/costas\_tree.c}, was written afresh for this study and shares no code
with the project's earlier solvers. It reproduces their exact tree sizes at every order where both have
been run, differing by exactly one because Paper~1 also counts the empty root:

\begin{table}[H]
\centering
\small
\begin{tabular}{@{}rrrr@{}}
\toprule
order & nodes, this search & nodes, Paper~1 & solutions \\
\midrule
11 & 789{,}511 & 789{,}512 & 2{,}509 \\
12 & 3{,}488{,}680 & 3{,}488{,}681 & 3{,}926 \\
13 & 18{,}290{,}043 & 18{,}290{,}044 & 7{,}077 \\
14 & 83{,}176{,}084 & 83{,}176{,}085 & 8{,}626 \\
\bottomrule
\end{tabular}
\caption{Exact tree sizes. At even orders the solution counts are exactly half the published totals
(7{,}852 and 17{,}252), as the symmetry restriction requires; at odd orders arrays whose first dot is in
the middle row are kept whole, so the counts exceed half.}
\label{tab:exact}
\end{table}

\section{Estimators and their validation}
\label{sec:validation}
Knuth's estimator~\cite{knuth1975} follows a random path from the root, choosing uniformly among the
valid children at each level. If the path meets $d_1,d_2,\dots$ children at successive levels, the
quantity $d_1+d_1d_2+d_1d_2d_3+\cdots$ is an unbiased estimate of the number of nodes, and averaging
independent probes reduces its variance. The \emph{stratified} variant enumerates every node at a fixed
depth $D$ exactly, counting the nodes above it exactly, and then probes beneath each frontier node,
which removes the variance contributed by the top of the tree. In both cases the standard error is
computed from the spread of the probes.

The estimators were validated against the exact counts at orders 12 to 16, three runs per order.
Every estimate is within $0.32\%$ of the exact count, and in all 15 runs the exact count lies within
two standard errors of the estimate.

\begin{table}[H]
\centering
\small
\begin{tabular}{@{}rrrr@{}}
\toprule
order & Knuth, $10^6$ probes & stratified, $D=3$ & stratified, $D=4$ \\
\midrule
12 & $-0.02\%$ ($\pm0.16\%$) & $+0.27\%$ ($\pm0.44\%$) & $-0.04\%$ ($\pm0.29\%$) \\
13 & $+0.03\%$ ($\pm0.17\%$) & $-0.13\%$ ($\pm0.40\%$) & $+0.04\%$ ($\pm0.26\%$) \\
14 & $-0.12\%$ ($\pm0.19\%$) & $+0.32\%$ ($\pm0.40\%$) & $+0.13\%$ ($\pm0.24\%$) \\
15 & $-0.15\%$ ($\pm0.20\%$) & $-0.09\%$ ($\pm0.37\%$) & $-0.20\%$ ($\pm0.22\%$) \\
16 & $-0.08\%$ ($\pm0.21\%$) & $-0.10\%$ ($\pm0.37\%$) & $+0.10\%$ ($\pm0.21\%$) \\
\bottomrule
\end{tabular}
\caption{Error of each estimate against the exact count, with two standard errors in parentheses. The
exact counts at orders 15 and 16 are those of Paper~1.}
\label{tab:validation}
\end{table}

Kilby, Slaney, Thi\'ebaux and Walsh describe weighted-backtrack and recursive estimators that improve
on Knuth's for uneven trees~\cite{kilby2006}. They were not implemented here. The Costas tree turns out
to be regular enough that plain probing is accurate and well calibrated, both at these orders and,
as Section~\ref{sec:cluster} shows, in genuine order-32 subtrees.

\section{The order-32 tree}
\label{sec:order32}
Five independent estimates at order 32 agree closely:

\begin{table}[H]
\centering
\small
\begin{tabular}{@{}lrr@{}}
\toprule
method & estimate & standard error \\
\midrule
Knuth, $2\times10^6$ probes, seed 21 & $1.5365\times10^{20}$ & $0.17\%$ \\
Knuth, $2\times10^6$ probes, seed 22 & $1.5409\times10^{20}$ & $0.17\%$ \\
Knuth, $2\times10^6$ probes, seed 23 & $1.5438\times10^{20}$ & $0.17\%$ \\
stratified, $D=3$, 14{,}640 frontier nodes & $1.5374\times10^{20}$ & $0.19\%$ \\
stratified, $D=4$, 408{,}470 frontier nodes & $1.5354\times10^{20}$ & $0.12\%$ \\
\bottomrule
\end{tabular}
\caption{Estimates of the order-32 tree. Their inverse-variance mean is $1.538\times10^{20}$ nodes, with a
relative standard error of $0.07\%$.}
\label{tab:order32}
\end{table}

\noindent There is no sign of a heavy tail: no single probe contributes more than $0.02\%$ of any total.
At orders 29, 30 and 31 the estimates are $1.566\times10^{18}$, $7.097\times10^{18}$ and
$3.406\times10^{19}$, so the tree grows by factors of $4.53$, $4.80$ and $4.52$ per order from 29 to 32,
and by $98.2$ in all. The alternation comes from the symmetry restriction, which keeps the middle row
whole at odd orders. The growth agrees with Paper~1's measured ratio of $4.80$ per order at orders 11 to
16, which extrapolated from order 16 gives $1.6\times10^{20}$.

\section{Measurements on a cluster node}
\label{sec:cluster}
The same program was run on one node of the cluster used for the campaign of Paper~8: two Intel Xeon
Sapphire Rapids processors, 104 cores in all, with the \texttt{powersave} frequency governor and turbo
enabled. The job cost 78 service units.

\begin{table}[H]
\centering
\small
\begin{tabular}{@{}lr@{}}
\toprule
measurement at order 32 & value \\
\midrule
node rate, one core alone & $3.72\times10^6$ nodes/s \\
node rate per core, all 104 cores busy & $3.14\times10^6$ nodes/s (range $3.13$--$3.16$) \\
node rate per node & $3.27\times10^8$ nodes/s \\
slowdown of a fully loaded node & $1.18$ \\
tree size, 8 independent Knuth runs & $1.5381\times10^{20}$ (relative s.e.\ $0.06\%$) \\
\bottomrule
\end{tabular}
\caption{Cluster measurements. The single-core rate equals the virtual machine's, and the tree estimate
agrees with Table~\ref{tab:order32} to within $0.01\%$.}
\label{tab:cluster}
\end{table}

The validation of Section~\ref{sec:validation} used orders 12 to 16. The cluster run also tested the
estimator on genuine order-32 subtrees. Each of 192 cores took a random path to depth 8 or 9 of the
order-32 tree, estimated the subtree below with 5{,}000 probes, and then counted it exactly.

\begin{table}[H]
\centering
\small
\begin{tabular}{@{}lrr@{}}
\toprule
& depth 8 & depth 9 \\
\midrule
subtrees counted exactly & 96 of 96 & 96 of 96 \\
median subtree size & $1.51\times10^9$ nodes & $7.69\times10^7$ nodes \\
exact count within two standard errors & $93.8\%$ & $94.8\%$ \\
total of estimates / total of exact counts & $0.9965$ & $1.0002$ \\
\bottomrule
\end{tabular}
\caption{The estimator on order-32 subtrees. No subtree reached the 30-minute time cap, so none was
excluded.}
\label{tab:deepval}
\end{table}

\noindent The coverage is close to the nominal $95\%$, and across the depth-8 subtrees, about
$1.5\times10^{11}$ nodes in all, the estimates' total is within $0.35\%$ of the exact one. The estimator is
therefore validated in the tree it is used on, not only at small orders.

\section{From nodes to cost}
\label{sec:cost}
A campaign keeps every core busy, so the relevant rate is the fully loaded one, $3.14\times10^6$ nodes per
core per second. At that rate the order-32 tree takes
\[
\frac{1.538\times10^{20}}{3.14\times10^6}\ \text{s}\approx1.36\times10^{10}\ \text{core-hours}
\approx1.5\times10^{6}\ \text{core-years}.
\]
On a queue charging 2 service units per core-hour, that is about $2.7\times10^{10}$ service units, or
27 million thousand-unit blocks (KSU). Exploiting the full symmetry group of the square, rather than the
single symmetry used here, would save at most a further factor of about four. Neither figure is within
reach of any allocation.

\section{Latency}
\label{sec:latency}
How long a campaign takes depends on how evenly its work divides. The cluster run estimated the subtree
below every one of the 408{,}470 nodes at depth 4, each with 20 probes, and scheduled them as tasks,
largest first, onto the least-loaded worker. The work is fairly even: the largest task is $34.7$ times
the mean, $1.15\times10^6$ core-hours, and the largest $1\%$ of tasks hold $3.8\%$ of the work.

\begin{table}[H]
\centering
\small
\begin{tabular}{@{}rrrr@{}}
\toprule
workers & wall time & efficiency & charged KSU \\
\midrule
104 & about 14{,}900 years & $100.0\%$ & $2.71\times10^7$ \\
1{,}040 & about 1{,}490 years & $100.0\%$ & $2.71\times10^7$ \\
10{,}400 & about 149 years & $99.7\%$ & $2.72\times10^7$ \\
104{,}000 & about 131 years & $11.3\%$ & $2.39\times10^8$ \\
\bottomrule
\end{tabular}
\caption{Wall time and charged cost of the straightforward search, with depth-4 tasks.}
\label{tab:latency}
\end{table}

\noindent Up to about $10^4$ workers, depth-4 tasks keep every core busy. Beyond that the largest task
exceeds the ideal wall time, and efficiency collapses unless the work is split more finely. The wall
times put the straightforward search beyond reach at any scale. A method costing about $10^4$ core-years,
as Section~\ref{sec:legacy} suggests for the 2011 method, would need about a year on $10^4$ cores,
provided its work divides as evenly.

\section{The 2011 enumeration was far more efficient}
\label{sec:legacy}
Paper~1 projected the order-32 search at $10^3$--$10^4$ modern core-years by scaling the reported cost of
the 2011 order-29 enumeration, 366.55 CPU-years~\cite{drakakis29}. The present search would need about
13{,}400 core-years at order 29, so the 2011 method was at least 36 times more efficient, and more once
the slower hardware of 2010 is allowed for. Paper~1's projection is therefore a statement about that
method, not about a straightforward column-by-column search. Section~\ref{sec:engineering} shows that
most of the gap is recoverable by engineering.

Scaling the 2011 cost by the measured growth from order 29 to 32, a factor of 98 rather than the $5^3=125$
of a five-fold-per-order rule, gives about 36{,}000 CPU-years on 2010 hardware. If per-core speed has
improved two- to four-fold since then, an assumption that has not been measured, that is
9{,}000--18{,}000 modern core-years, or about 160{,}000--320{,}000 KSU on a 2-unit queue, and up to about 370{,}000 allowing for the
slowdown of a fully loaded node measured in Section~\ref{sec:cluster}. This assumes that
the 2011 method's tree grows like the present one; its algorithm has not been reproduced here.

\section{Engineering the search}
\label{sec:engineering}
Most of the 2011 method's advantage turns out to be recoverable by engineering the same search. Three
changes were tested, each first validated against exact counts.

\paragraph{Bit-parallel candidate rows.} The search of Section~\ref{sec:tree} tests each candidate row
against every earlier column, $O(n\cdot c)$ work to expand a node at column $c$. Row $r$ conflicts with
column $i$ exactly when the difference $r-p_i$ is already used at lag $c-i$, so the rows forbidden by
column $i$ are the used-difference mask of that lag, shifted into row coordinates. ORing these masks gives
every forbidden row at once, in $O(c)$ word operations per node. The tree is unchanged: node and
solution counts are identical at orders 10 to 14. At order 32 the rate rises from $3.33\times10^6$ to
$1.68\times10^7$ nodes per second on the same core, a factor of $5.0$.

\paragraph{Column order.} Four orders were compared, all with the same symmetry restriction, which applies
to whichever column is placed first.

\begin{table}[H]
\centering
\small
\begin{tabular}{@{}rrrrr@{}}
\toprule
order & left to right & outside in & middle out & fewest legal rows \\
\midrule
12 & $3.49\times10^6$ ($0.36$\,s) & $1.22\times10^7$ ($1.39$\,s) & $3.49\times10^6$ ($0.44$\,s) & $2.42\times10^6$ ($1.22$\,s) \\
13 & $1.83\times10^7$ ($2.19$\,s) & $7.81\times10^7$ ($8.86$\,s) & $1.79\times10^7$ ($1.95$\,s) & $1.22\times10^7$ ($7.58$\,s) \\
14 & $8.32\times10^7$ ($12.5$\,s) & $4.38\times10^8$ ($55.9$\,s) & $8.32\times10^7$ ($11.5$\,s) & $5.38\times10^7$ ($35.6$\,s) \\
\bottomrule
\end{tabular}
\caption{Nodes and time for four column orders. ``Fewest legal rows'' chooses at each node the unplaced
column with fewest legal rows and abandons a branch when any column has none.}
\label{tab:orders}
\end{table}

\noindent Any order that keeps the placed columns contiguous gives the same tree, because whether a block
of dots is Costas does not depend on where the block sits; at even orders the counts for left to right and
middle out are identical. Placing dots far apart, outside in, wastes the short-lag constraints, which prune
hardest, and gives a tree five times larger by order 14. Choosing the column with fewest legal rows shrinks
the tree increasingly, to $65\%$ of the baseline at order 14, and its growth per order is about $4.65$
against $4.84$. But each node costs about three times as much, so it is slower at every order measured.

Orders that place a column between two earlier ones need one more check than left to right. A new dot
forms two pairs at the same lag $L$, one on each side, and these collide exactly when its row is the
midpoint of the other two, $2r=p_{c-L}+p_{c+L}$: the Ladendorf condition of the 2009 search. A first
version of these searches omitted the check and accepted non-Costas arrays; the validation against exact
array counts caught it before any result was used.

\paragraph{Full symmetry pruning.} Using all eight symmetries of the square, the search keeps exactly one
array per class: the lexicographically smallest of its eight images. A prefix is cut when some image is
guaranteed to be smaller for every completion. An image entry is either known, or confined to an interval:
an unplaced row's column lies among the unplaced columns. The vertical flip compares known values; the four
symmetries involving transposition can often be decided early from the partly known inverse; the horizontal
flip and the half-turn are decided only near the leaves. The pruning is exact: at every order from 4 to
13 the search outputs precisely the class representatives found by enumerating every Costas array and
grouping them by symmetry, including 990 classes among 7{,}852 arrays at order 12 and 1{,}616 among 12{,}828
at order 13. On the pruned tree the estimator is within $0.3\%$ of the exact counts at orders 13 to 15.

\begin{table}[H]
\centering
\small
\begin{tabular}{@{}rrrr@{}}
\toprule
order & one symmetry & all eight & fewer nodes \\
\midrule
12 & $3{,}488{,}680$ & $1{,}544{,}404$ & $2.26\times$ \\
13 & $18{,}290{,}043$ & $7{,}509{,}524$ & $2.44\times$ \\
14 & $83{,}176{,}084$ & $37{,}132{,}082$ & $2.24\times$ \\
15 & $435{,}090{,}280$ & $182{,}409{,}077$ & $2.39\times$ \\
\bottomrule
\end{tabular}
\caption{Nodes with one and with all eight symmetries, in the bit-parallel search.}
\label{tab:symmetry}
\end{table}

\noindent At order 32 three independent estimates put the pruned tree at $7.22\times10^{19}$ nodes, $2.13$
times smaller, while the extra comparisons lower the rate to $1.42\times10^7$ nodes per second. The search is
about $1.8$ times faster overall: less than the theoretical four, because two of the symmetries can be
decided only near the leaves.

\paragraph{The revised cost.} With the fully loaded slowdown of Section~\ref{sec:cluster}:

\begin{table}[H]
\centering
\small
\begin{tabular}{@{}lrrr@{}}
\toprule
search & core-hours & core-years & KSU \\
\midrule
as in Section~\ref{sec:cost} & $1.36\times10^{10}$ & $1.55\times10^6$ & $2.7\times10^7$ \\
bit-parallel candidate rows & $3.0\times10^{9}$ & $3.4\times10^5$ & $6.0\times10^6$ \\
and full symmetry pruning & $1.7\times10^{9}$ & $1.9\times10^5$ & $3.3\times10^6$ \\
\bottomrule
\end{tabular}
\caption{The cost of the order-32 search, before and after engineering.}
\label{tab:revised}
\end{table}

\noindent The engineered search is about eight times cheaper, but still about 18 years on $10^4$ cores. At
order 29 it would need about 1{,}650 core-years against the 2011 enumeration's 366.55 CPU-years, so a factor
of four to five remains, more if 2010 hardware was slower per core. Whether that factor lies in node rate or
in a smaller tree is not known. For comparison, a 2026 preprint estimates by sampling that an
unrestricted order-28 search with its software would take 43.7 CPU-years, and projects 7.97 GPU-years
at order 30~\cite{dumitru2026localpatterns}. Reductions in the growth per order compound: the fewest-legal-rows order
lowers it from $4.84$ to about $4.65$, which over the orders up to 32 is worth about three times, if its
per-node cost could be brought down to the baseline's. Paper~10 measures a related attempt: binary
displacement consistency, which maintains the domains that order needs, reduced the forward-checking tree
about fourfold at order 14 but made complete searches slower, and was not retained.

\section{Implications for a distributed SAT campaign}
\label{sec:sat}
A planning model for a distributed SAT campaign can be written
\[
K_{\mathrm{SAT}}=K_{\mathrm{ref}}\,\frac{\mu}{S},
\]
where $K_{\mathrm{ref}}$ is the cost of a reference enumeration, $S$ the improvement of a modern direct
method over it, and $\mu$ the overhead of SAT and its decomposition over that direct method. Both
factors can now be bounded from measurements in this volume.
\begin{itemize}
\item Relative to the 2011 method, the straightforward modern search is not faster but at least 36
times slower, so $S>1$ requires an algorithm better than the 2011 one, which no one in this
investigation has shown.
\item For the transpose-symmetric subclass, the SAT campaign of Paper~8 used about 490 times the
computation of the direct search later run to completion, so $\mu$ was of that order, not a small constant. Even $\mu=5$ applied to
the 2011 method's scaled cost gives about 0.8--1.6 million KSU.
\end{itemize}
On cost, therefore, a distributed SAT campaign is not a viable route to the unrestricted problem. The
only candidate is a direct enumeration at least as efficient as the 2011 method, distributed through a
fail-closed task queue, at a cost of order $10^5$ KSU on present assumptions. Before any allocation is
requested, that method should be reproduced and its rate and tree measured at orders where they can be
checked, such as 20 to 24; the estimator validated here makes that calibration cheap.

\section{Limitations}
\label{sec:limits}
The main estimates describe one search, with columns filled from left to right and one symmetry used;
Section~\ref{sec:engineering} measures three changes to it. Other
filling orders, such as from the middle outwards, give different trees, and the 2011 method's tree is
not measured here. The node rates were measured on one virtual machine and one cluster node. The weighted-backtrack and recursive estimators were not implemented. The improvement in
per-core speed since 2010 is an assumption.
